\chapter{Código para  reproducir los ejemplos de la sección Resultados \label{sec:ejemplos}}
\section{Ejemplo  de uso de variables de usuario para controlar la proporcion de cada genotipo en la población total}
Código para reproducir el segundo ejemplo del punto 5.3.1:
\begin{R}
dfuv2 <- data.frame(Genotype = c("WT", "B", "A", "B, A", "C, A"),
                       Fitness = c("0*n_",
                                    "1.5",
                                    "1.002",
                                    "1.003",
                                    "1.004"))

afuv2 <- allFitnessEffects(genotFitness = dfuv2,
                          frequencyDependentFitness = TRUE,
                          frequencyType = "abs")

userVars <- list(
    list(Name           = "genAProp",
        Value       = 0.5
    ),list(Name           = "genBProp",
        Value       = 0.5
    ),list(Name           = "genABProp",
        Value       = 0.0
    ),list(Name           = "genACProp",
        Value       = 0.0
    )
)

userVars <- createUserVars(userVars)

rules <- list(
        list(ID = "rule_1",
            Condition = "TRUE",
            Action = "genBProp = n_B/N"
        ),list(ID = "rule_2",
            Condition = "TRUE",
            Action = "genAProp = n_A/N"
        ),list(ID = "rule_3",
            Condition = "TRUE",
            Action = "genABProp = n_A_B/N"
        ),list(ID = "rule_4",
            Condition = "TRUE",
            Action = "genACProp = n_A_C/N"
        )
    )

rules <- createRules(rules, afuv2)

set.seed(1)
uvex2 <- oncoSimulIndiv(
                    afuv2, 
                    model = "McFLD",
                    mu = 1e-4,
                    sampleEvery = 0.01,
                    initSize = c(20000, 20000),
                    initMutant = c("A", "B"),
                    finalTime = 10,
                    onlyCancer = FALSE,
                    userVars = userVars,
                    rules = rules
                    )

plot(uvex2, show = "genotypes", type = "line")      

plot(unlist(uvex2$other$userVarValues) [c(FALSE, FALSE, FALSE, FALSE, TRUE)], unlist(uvex2$other$userVarValues) [c(TRUE, FALSE, FALSE, FALSE, FALSE)], xlab="Time", ylab="Proportion", ylim=c(0,1), type="l", col="purple")  
lines(unlist(uvex2$other$userVarValues) [c(FALSE, FALSE, FALSE, FALSE, TRUE)], unlist(uvex2$other$userVarValues) [c(FALSE, TRUE, FALSE, FALSE, FALSE)], type="l", col="#E6AB02") 
lines(unlist(uvex2$other$userVarValues) [c(FALSE, FALSE, FALSE, FALSE, TRUE)], unlist(uvex2$other$userVarValues) [c(FALSE, FALSE, TRUE, FALSE, FALSE)], type="l", col="#1B9E77") 
lines(unlist(uvex2$other$userVarValues) [c(FALSE, FALSE, FALSE, FALSE, TRUE)], unlist(uvex2$other$userVarValues) [c(FALSE, FALSE, FALSE, TRUE, FALSE)], type="l", col="#666666") 
legend(0,1,legend=c("genABProp", "genACProp", "genAProp", "genBProp"), col=c("purple", "#E6AB02", "#1B9E77", "#666666"), lty= 1:2) 
\end{R}
\section{Ejemplo  de uso de variables de usuario para calcular la diferencia entre tasa de nacimiento y muerte de los genotipos}
Código para reproducir el segundo ejemplo del punto 5.3.1:
\begin{R}
dfuv3 <- data.frame(Genotype = c("WT", "A", "B"), 
                   Fitness = c("1",
                               "1 + 0.2 * (n_B > 10)",
                               ".9 + 0.4 * (n_A > 10)"
                               ))
afuv3 <- allFitnessEffects(genotFitness = dfuv3, 
                         frequencyDependentFitness = TRUE)

userVars <- list(
    list(Name           = "genWTRateDiff",
        Value       = 0.5
    ),list(Name           = "genARateDiff",
        Value       = 0.5
    ),list(Name           = "genBRateDiff",
        Value       = 0.0
    )
)

userVars <- createUserVars(userVars)

rules <- list(
        list(ID = "rule_1",
            Condition = "TRUE",
            Action = "genWTRateDiff = b_-d_"
        ),list(ID = "rule_2",
            Condition = "TRUE",
            Action = "genARateDiff = b_1-d_1"
        ),list(ID = "rule_3",
            Condition = "TRUE",
            Action = "genBRateDiff = b_2-d_2"
        )
    )

rules <- createRules(rules, afuv3)

set.seed(1)

uvex3 <- oncoSimulIndiv(afuv3,
                       model = "McFLD", 
                       onlyCancer = FALSE, 
                       finalTime = 150,
                       mu = 1e-4,
                       initSize = 5000, 
                       keepPhylog = FALSE,
                       seed = NULL, 
                       errorHitMaxTries = FALSE, 
                       errorHitWallTime = FALSE,
                       userVars = userVars,
                       rules = rules)

plot(uvex3, show = "genotypes", type = "line")    

plot(unlist(uvex3$other$userVarValues) [c(FALSE, FALSE, FALSE, TRUE)], unlist(uvex3$other$userVarValues) [c(FALSE, FALSE, TRUE, FALSE)], xlab="Time", ylab="Rate Diff", xlim=c(0,150), ylim=c(-0.75,0.75), type="l", col="#1B9E77")  
lines(unlist(uvex3$other$userVarValues) [c(FALSE, FALSE, FALSE, TRUE)], unlist(uvex3$other$userVarValues) [c(TRUE, FALSE, FALSE, FALSE)], type="l", col="#A6761D") 
lines(unlist(uvex3$other$userVarValues) [c(FALSE, FALSE, FALSE, TRUE)], unlist(uvex3$other$userVarValues) [c(FALSE, TRUE, FALSE, FALSE)], type="l", col="#666666") 
legend(0,0.75,legend=c("genWTRateDiff", "genARateDiff", "genBRateDiff"), col=c("#1B9E77", "#A6761D", "#666666"), lty= 1:2) 
\end{R}
\section{Ejemplo de terapia adaptativa para atacar al genotipo más adaptado}
Código para reproducir el primer ejemplo del punto 5.3.2:
\begin{R}
dfat2 <- data.frame(Genotype = c("WT", "A", "B"), 
                   Fitness = c("1",
                               "1 + 0.2 * (n_B > 10)",
                               ".9 + 0.4 * (n_A > 10)"
                               ))
afat2 <- allFitnessEffects(genotFitness = dfat2, 
                         frequencyDependentFitness = TRUE)
    

userVars <- list(
    list(Name           = "genADiff",
        Value       = 0
    ),list(Name           = "genBDiff",
        Value       = 0
    ),list(Name           = "genWTDiff",
        Value       = 0
    )
)

userVars <- createUserVars(userVars)

rules <- list(
        list(ID = "rule_1",
            Condition = "TRUE",
            Action = "genADiff = b_1-d_1"
        ),list(ID = "rule_2",
            Condition = "TRUE",
            Action = "genBDiff = b_2-d_2"
        ),list(ID = "rule_3",
            Condition = "TRUE",
            Action = "genWTDiff = b_-d_"
        ),list(ID = "rule_4",
            Condition = "n_A < 1000",
            Action = "genADiff = -1"
        ),list(ID = "rule_5",
            Condition = "n_B < 1000",
            Action = "genBDiff = -1"
        ),list(ID = "rule_6",
            Condition = "n_ < 1000",
            Action = "genWTDiff = -1"
        )
    )

rules <- createRules(rules, afat2)

interventions <- list(
    list(ID           = "i1",
        Trigger       = "genADiff > genBDiff and genADiff > genWTDiff",
        WhatHappens   = "n_A = n_A*0.5",
        Periodicity   = 10,
        Repetitions   = Inf
    ),list(ID           = "i2",
        Trigger       = "genBDiff > genADiff and genBDiff > genWTDiff",
        WhatHappens   = "n_B = n_B*0.5",
        Periodicity   = 10,
        Repetitions   = Inf
    ),list(ID           = "i3",
        Trigger       = "genWTDiff > genADiff and genWTDiff > genBDiff",
        WhatHappens   = "n_ = n_*0.5",
        Periodicity   = 10,
        Repetitions   = Inf
    )
)

interventions <- createInterventions(interventions, afat2)

set.seed(1) ## for reproducibility
atex2 <- oncoSimulIndiv(afat2,
                       model = "McFLD", 
                       onlyCancer = FALSE, 
                       finalTime = 200,
                       mu = 1e-4,
                       initSize = 5000, 
                       keepPhylog = FALSE,
                       seed = NULL, 
                       errorHitMaxTries = FALSE, 
                       errorHitWallTime = FALSE,
                       userVars = userVars,
                       rules = rules,
                       interventions = interventions)

plot(atex2, show = "genotypes", type = "line")

plot(unlist(atex2$other$userVarValues) [c(FALSE, FALSE, FALSE, TRUE)], unlist(atex2$other$userVarValues) [c(TRUE, FALSE, FALSE, FALSE)], xlab="Time", ylab="RateDiff", ylim=c(-0.75,0.75), type="l", col="#A6761D")  
lines(unlist(atex2$other$userVarValues) [c(FALSE, FALSE, FALSE, TRUE)], unlist(atex2$other$userVarValues) [c(FALSE, TRUE, FALSE, FALSE)], type="l", col="#666666") 
lines(unlist(atex2$other$userVarValues) [c(FALSE, FALSE, FALSE, TRUE)], unlist(atex2$other$userVarValues) [c(FALSE, FALSE, TRUE, FALSE)], type="l", col="#1B9E77") 
legend(0,0.75,legend=c("genADiff", "genBDiff", "genWTDiff"), col=c("#A6761D", "#666666", "#1B9E77"), lty= 1:2) 
\end{R}
\section{Ejemplo de terapia adaptativa en escenario verosímil}
Código para reproducir el segundo ejemplo del punto 5.3.2:
\begin{R}
dfat3 <- data.frame(Genotype = c("WT", "A", "B"), 
                   Fitness = c("1",
                               "0.8 + 0.2 * (n_B > 10) + 0.1 (n_A > 10)",
                               "0.8 + 0.25 * (n_B > 10)"
                               ))
afat3 <- allFitnessEffects(genotFitness = dfat3, 
                         frequencyDependentFitness = TRUE)
    

userVars <- list(
    list(Name           = "lastMeasuredA",
        Value       = 0
    ),list(Name           = "lastMeasuredB",
        Value       = 0
    ),list(Name           = "previousA",
        Value       = 0
    ),list(Name           = "previousB",
        Value       = 0
    ),list(Name           = "lastTime",
        Value       = 0
    ),list(Name           = "measure",
        Value       = 0
    ),list(Name           = "treatment",
        Value       = 0
    )
)

userVars <- createUserVars(userVars)

rules <- list(
       list(ID = "rule_1",
            Condition = "T - lastTime < 10",
            Action = "measure = 0"
        ),list(ID = "rule_2",
            Condition = "T - lastTime >= 10",
            Action = "measure = 1;lastTime = T"
        ),list(ID = "rule_3",
            Condition = "measure == 1",
            Action = "previousA = lastMeasuredA;previousB = lastMeasuredB;lastMeasuredA = n_A;lastMeasuredB = n_B"
        ), list(ID = "rule_4",
            Condition = "TRUE",
            Action = "treatment = 0"
        ), list(ID = "rule_5",
            Condition = "lastMeasuredA + lastMeasuredB > 100",
            Action = "treatment = 1"
        ), list(ID = "rule_6",
            Condition = "lastMeasuredA - PreviousA > 500",
            Action = "treatment = 2"
        ),list(ID = "rule_7",
            Condition = "lastMeasuredB - PreviousB > 500",
            Action = "treatment = 3"
        ), list(ID = "rule_8",
            Condition = "lastMeasuredA - PreviousA > 500 and lastMeasuredB - PreviousB > 500",
            Action = "treatment = 4"
        )
    )

rules <- createRules(rules, afat3)


interventions <- list(
    list(ID           = "basicTreatment",
        Trigger       = "treatment == 1",
        WhatHappens   = "N = 0.8*N",
        Periodicity   = 10,
        Repetitions   = Inf
    ),list(ID           = "treatmentOverA",
        Trigger       = "treatment == 2 or treatment == 4",
        WhatHappens   = "n_B = n_B*0.3",
        Periodicity   = 20,
        Repetitions   = Inf
    ),list(ID           = "treatmentOverB",
        Trigger       = "treatment == 3 or treatment == 4",
        WhatHappens   = "n_B = n_B*0.3",
        Periodicity   = 20,
        Repetitions   = Inf
    ),list(ID           = "intervention",
        Trigger       = "lastMeasuredA+lastMeasuredB > 5000",
        WhatHappens   = "N = 0.1*N",
        Periodicity   = 70,
        Repetitions   = Inf
    )
)

interventions <- createInterventions(interventions, afat3)

set.seed(1) ## for reproducibility
atex2 <- oncoSimulIndiv(afat3,
                       model = "McFLD", 
                       onlyCancer = FALSE, 
                       finalTime = 200,
                       mu = 1e-4,
                       initSize = 5000, 
                       keepPhylog = FALSE,
                       seed = NULL, 
                       errorHitMaxTries = FALSE, 
                       errorHitWallTime = FALSE,
                       userVars = userVars,
                       rules = rules,
                       interventions = interventions)

plot(atex2, show = "genotypes", type = "line")
\end{R}

